% English translation of chapter1.tex lines 1230--1582.
% Source: Methods of Algebra, Volume 2, commit
% 9a5803ff77dd3257484cb177f851a73770a59dd3 (CC BY 4.0).
% stable-unit-id: o014.aljabr2.chapter1.gabriel-zisman-localization

% segment-id: o014.li2.chapter1.gabriel-zisman-localization.h001
\section{Gabriel--Zisman Localization}\label{sec:cat-localization}
\hypertarget{o014.aljabr2.chapter1.gabriel-zisman-localization}{ }

% segment-id: o014.li2.chapter1.gabriel-zisman-localization.p001
Gabriel--Zisman localization, henceforth simply called localization, is the
most economical way to adjoin inverses to a family of morphisms in a category.
The theory originates in \cite{GZ67} and is characterized by a universal
property.

% segment-id: o014.li2.chapter1.gabriel-zisman-localization.d001
\begin{definition}[P.\ Gabriel, M.\ Zisman]\label{def:cat-localization}
	\index{localization}
	\index[sym1]{CS@$\mathcal{C}[S^{-1}]$}
	Let $\mathcal{C}$ be a category and let $S$ be a subset of
	$\Mor(\mathcal{C})$ that contains all identity morphisms and is closed
	under composition. A \emph{localization} of $\mathcal{C}$ at $S$ means a
	category $\mathcal{C}[S^{-1}]$ (which may be a ``large'' category), together
	with a functor $Q:\mathcal{C}\to\mathcal{C}[S^{-1}]$, called the
	localization functor, satisfying the following conditions.
	% segment-id: o014.li2.chapter1.gabriel-zisman-localization.l001
	\begin{itemize}
		% segment-id: o014.li2.chapter1.gabriel-zisman-localization.i001
		\item For every $s\in S$, its image $Q(s)$ is an isomorphism in
		$\mathcal{C}[S^{-1}]$.
		% segment-id: o014.li2.chapter1.gabriel-zisman-localization.i002
		\item For every category $\mathcal{D}$ (again possibly ``large'') and
		functor $F:\mathcal{C}\to\mathcal{D}$, if $F$ sends every morphism in
		$S$ to an isomorphism, then there is a unique functor
		$F[S^{-1}]:\mathcal{C}[S^{-1}]\to\mathcal{D}$ such that
		$F=F[S^{-1}]Q$.
	\end{itemize}
\end{definition}
\index{size issues}

% segment-id: o014.li2.chapter1.gabriel-zisman-localization.p002
By convention, we often omit $Q$ from the data. Some sources, such as
\cite[Definition 7.1.1]{KS06}, give a somewhat weaker universal property for
localization\footnote{From the perspective of $2$-categories, the more natural
universal property is this: for every $\mathcal{D}$, the functor
$Q^*:\mathcal{D}^{\mathcal{C}[S^{-1}]}\to\mathcal{D}^{\mathcal{C}}$ is fully
faithful, and $G\in\Obj(\mathcal{D}^{\mathcal{C}})$ is isomorphic to some
$Q^*(F)=FQ$ if and only if $G$ sends the elements of $S$ to isomorphisms. The
functor $Q$ in Definition \ref{def:cat-localization} automatically has this
property---the fully faithful part is Proposition
\ref{prop:localization-extra-property}, and the rest is easy.}.

% segment-id: o014.li2.chapter1.gabriel-zisman-localization.p003
We next show that localization is unique up to a unique isomorphism.

% segment-id: o014.li2.chapter1.gabriel-zisman-localization.pr001
\begin{proposition}\label{prop:cat-localization-unique}
	If $(\mathcal{C}[S^{-1}],Q)$ and $(\mathcal{C}[S^{-1}]',Q')$ are both
	localizations of $\mathcal{C}$ at $S$, then there is a unique pair of
	functors
	% segment-id: o014.li2.chapter1.gabriel-zisman-localization.g001
	$\begin{tikzcd}
		\mathcal{C}[S^{-1}] \arrow[r, shift left, "G"] & \mathcal{C}[S^{-1}]' \arrow[l, shift left, "{G'}"]
	\end{tikzcd}$
	such that $GQ=Q'$, $G'Q'=Q$, and
	$G'G=\identity_{\mathcal{C}[S^{-1}]}$,
	$GG'=\identity_{\mathcal{C}[S^{-1}]'}$.
\end{proposition}
% segment-id: o014.li2.chapter1.gabriel-zisman-localization.pf001
\begin{proof}
	Since $Q'$ sends $S$ to isomorphisms, the universal property gives a unique
	$G$ such that $GQ=Q'$. Similarly, there is a unique $G'$ such that
	$G'Q'=Q$. Since $G'GQ=Q$, taking
	$\mathcal{D}=\mathcal{C}[S^{-1}]$ and $F=Q$ in the universal property gives
	$G'G=\identity_{\mathcal{C}[S^{-1}]}$. Likewise,
	$GG'=\identity_{\mathcal{C}[S^{-1}]'}$\footnote{Translator's note: the
	source has $\identity_{\mathcal{C}[S^{-1}]}$ for this last identity. Since
	$GG'$ is an endofunctor of $\mathcal{C}[S^{-1}]'$, the prime has been added
	to the target localization here (O014-C017).}.
\end{proof}

% segment-id: o014.li2.chapter1.gabriel-zisman-localization.pr002
\begin{proposition}
	Suppose that the localization $(\mathcal{C}[S^{-1}],Q)$ exists. Let
	$S^{\opp}$ denote the corresponding image of $S$ in
	$\Mor(\mathcal{C}^{\opp})$. Then there is an equivalence of categories,
	compatible with $Q$,
	$\mathcal{C}[S^{-1}]^{\opp}\leftrightarrows
	\mathcal{C}^{\opp}[(S^{\opp})^{-1}]$.
\end{proposition}
% segment-id: o014.li2.chapter1.gabriel-zisman-localization.pf002
\begin{proof}
	By the uniqueness of localization, it suffices to apply $(-)^{\opp}$ to all
	the categories and functors in Definition \ref{def:cat-localization}, since
	this operation does not reverse the direction of a functor.
\end{proof}

% segment-id: o014.li2.chapter1.gabriel-zisman-localization.p004
The problem is thus reduced to the existence and construction of localization.
For general $\mathcal{C}$ and $S$, one can construct $\mathcal{C}[S^{-1}]$ by
formally adjoining inverses to $\mathcal{C}$, much as in the construction of a
free group. More concretely,
$\Obj(\mathcal{C}[S^{-1}])=\Obj(\mathcal{C})$, while morphisms in
$\mathcal{C}[S^{-1}]$ are represented formally by finite zigzags
% segment-id: o014.li2.chapter1.gabriel-zisman-localization.g002
\[ \cdots b t^{-1} a s^{-1} \cdots = \left(
	\begin{tikzcd}[row sep=tiny, column sep=tiny]
	\cdots \arrow[rd] & & Y \arrow[ld, "s"'] \arrow[rd, "a"] & & W \arrow[ld, "t"'] \arrow[rd, "b"] & & \cdots \arrow[ld] \\
	& X & & Z & & U &
\end{tikzcd}\right) \]
with $a,b,\ldots\in\Mor(\mathcal{C})$ and $s,t,\ldots\in S$. Composition
and $Q$ are defined in the evident way, but one must quotient by the
equivalence relation generated by
% segment-id: o014.li2.chapter1.gabriel-zisman-localization.q001
\[ s^{-1}t^{-1}=(ts)^{-1}, \quad ss^{-1}=\identity, \quad s^{-1}s=\identity \]
and so on; see the sketch in \cite[I.1.1]{GZ67}. Localization may therefore be
viewed as a kind of calculus of fractions for morphisms.

% segment-id: o014.li2.chapter1.gabriel-zisman-localization.p005
This construction is completely general, but its drawbacks are equally clear,
chiefly because the equivalence relation is hard to work with. For example, it
is difficult to characterize those $f,g\in\Mor(\mathcal{C})$ for which
$Qf=Qg$. Fortunately, in practice $S$ is often a left (or right)
multiplicative system, to be introduced in Definition
\ref{def:multiplicative-family}. In that case the description of
$\mathcal{C}[S^{-1}]$ can be simplified: the zigzags above can be
\emph{straightened}, so that every morphism can be written in the form
$as^{-1}$ (or $s^{-1}a$).

% segment-id: o014.li2.chapter1.gabriel-zisman-localization.p006
Before introducing multiplicative systems, we record a simple property of the
construction above.

% segment-id: o014.li2.chapter1.gabriel-zisman-localization.pr003
\begin{proposition}\label{prop:localization-extra-property}
	The localization functor $Q:\mathcal{C}\to\mathcal{C}[S^{-1}]$ is
	essentially surjective. Moreover, for every category $\mathcal{D}$ and
	functors $A,B:\mathcal{C}[S^{-1}]\to\mathcal{D}$, the evident map
	% segment-id: o014.li2.chapter1.gabriel-zisman-localization.q002
	\[ \Hom_{\mathcal{D}^{\mathcal{C}[S^{-1}]}}(A,B) \to \Hom_{\mathcal{D}^{\mathcal{C}}}\left(AQ,BQ\right) \]
	is a bijection. In other words,
	$Q^*:\mathcal{D}^{\mathcal{C}[S^{-1}]}\to\mathcal{D}^{\mathcal{C}}$ is a
	fully faithful functor.
\end{proposition}
% segment-id: o014.li2.chapter1.gabriel-zisman-localization.pf003
\begin{proof}
	We may realize the localization concretely as above. Write the map induced
	by $Q$ on objects as $X\mapsto\underline{X}$. It is a bijection, hence $Q$
	is essentially surjective. It follows at once that the map
	$\Hom_{\mathcal{D}^{\mathcal{C}[S^{-1}]}}(A,B)\to
	\Hom_{\mathcal{D}^{\mathcal{C}}}(AQ,BQ)$ is injective.

	For surjectivity, let
	$\varphi=(\varphi_X)_{X\in\Obj(\mathcal{C})}\in
	\Hom_{\mathcal{D}^{\mathcal{C}}}(AQ,BQ)$. For each
	$\underline{X}=QX$, define a morphism
	$\tilde{\varphi}_{\underline{X}}:A\underline{X}\to B\underline{X}$ in
	$\mathcal{D}$ to be $\varphi_X$. It remains to show that
	$\tilde{\varphi}:=(\tilde{\varphi}_{\underline{X}})_{\underline{X}}$ lies
	in $\Hom_{\mathcal{D}^{\mathcal{C}[S^{-1}]}}(A,B)$. In other words, for
	every morphism $f:\underline{X}\to\underline{Y}$ in
	$\mathcal{C}[S^{-1}]$, we must verify that the following diagram commutes in
	$\mathcal{D}$:
	% segment-id: o014.li2.chapter1.gabriel-zisman-localization.g003
	\[\begin{tikzcd}
		A\underline{X} \arrow[r, "{\tilde{\varphi}_{\underline{X}}}"] \arrow[d, "Af"'] & B\underline{X} \arrow[d, "Bf"] \\
		A\underline{Y} \arrow[r, "{\tilde{\varphi}_{\underline{Y}}}"'] & B\underline{Y} .
	\end{tikzcd}\]
	By the zigzag construction of morphisms, $f$ decomposes into a string of
	morphisms of the form $Qa$ or $(Qs)^{-1}$, with
	$a\in\Mor(\mathcal{C})$ and $s\in S$. The commutativity of the diagram
	therefore reduces to the corresponding property of $\varphi$.
\end{proof}

% segment-id: o014.li2.chapter1.gabriel-zisman-localization.d002
\begin{definition}\label{def:multiplicative-family}
	\index{multiplicative system}
	A subset $S\subset\Mor(\mathcal{C})$ satisfying the following properties
	is called a \emph{left multiplicative system} in $\mathcal{C}$.
	% segment-id: o014.li2.chapter1.gabriel-zisman-localization.l002
	\begin{enumerate}[(S1)]
		% segment-id: o014.li2.chapter1.gabriel-zisman-localization.i003
		\item For every $X\in\Obj(\mathcal{C})$, one has $\identity_X\in S$.
		% segment-id: o014.li2.chapter1.gabriel-zisman-localization.i004
		\item For all $f,g\in S$, if the composite $gf$ is defined, then
		$gf\in S$.
		% segment-id: o014.li2.chapter1.gabriel-zisman-localization.i005
		\item Given morphisms
		$X\xrightarrow{s\in S}Z\xleftarrow{f}Y$, there are morphisms in
		$\mathcal{C}$, $X\xleftarrow{f'}W\xrightarrow{s'\in S}Y$, such that
		$sf'=fs'$, depicted by the commutative diagram
		% segment-id: o014.li2.chapter1.gabriel-zisman-localization.g004
		\[\begin{tikzcd}
			X \arrow[d, "{s \in S}"'] & W \arrow[dashed, l, "{f'}"'] \arrow[dashed, d, "{s' \in S}"] \\
			Z & Y \arrow[l, "f"]
		\end{tikzcd} \quad \text{commutative}. \]
		% segment-id: o014.li2.chapter1.gabriel-zisman-localization.i006
		\item Let $f,g:X\to Y$ be morphisms in $\mathcal{C}$ and let
		$s:Y\to W$ belong to $S$. If $sf=sg$, then there is a morphism
		$t:Z\to X$ in $S$ such that $ft=gt$. Diagrammatically,
		% segment-id: o014.li2.chapter1.gabriel-zisman-localization.g005
		\[\begin{tikzcd}
			Z \arrow[dashed, r, "t \in S"] \arrow[dash, rr, "\text{commutative}"', to path={ -- ([yshift=-2ex]\tikztostart.south) -- ([yshift=-2ex]\tikztotarget.south) [midway]\tikztonodes -- (\tikztotarget)}] & X \arrow[r, shift left, "f"] \arrow[r, shift right, "g"'] & Y \arrow[r, "s \in S"] & W .
		\end{tikzcd}\]
	\end{enumerate}
	The dashed parts of these diagrams are the morphisms whose existence is
	asserted.

	If the image $S^{\opp}$ of $S$ in $\Mor(\mathcal{C}^{\opp})$ is a left
	multiplicative system, then $S$ is called a \emph{right multiplicative
	system} in $\mathcal{C}$. Equivalently, reverse all arrows in conditions
	(S3) and (S4). A subset $S$ that is both a left and a right multiplicative
	system is called a \emph{multiplicative system} in $\mathcal{C}$.
\end{definition}

% segment-id: o014.li2.chapter1.gabriel-zisman-localization.c001
\begin{convention}
	For a given left or right multiplicative system $S$, from now on arrows
	belonging to $S$ will always be drawn as $\rightarrowtail$ in commutative
	diagrams, to distinguish them from the other arrows.
\end{convention}

% segment-id: o014.li2.chapter1.gabriel-zisman-localization.p007
Fix $X\in\Obj(\mathcal{C})$. In accordance with this convention, define the
category $S_{/X}$ (or $S_{X/}$) for a left (or right) multiplicative system
$S$ as follows.
\index[sym1]{SX@$S_{/X}, S_{X/}$}
% segment-id: o014.li2.chapter1.gabriel-zisman-localization.q003
\begin{align*}
	\Obj\left(S_{/X}\right) &:= \left\{\text{morphisms}\;X\leftarrowtail Z\right\}, & \Mor(S_{/X}) &:= \left\{\text{commutative diagrams}\;
	% segment-id: o014.li2.chapter1.gabriel-zisman-localization.g006
	{\begin{tikzcd}[row sep=small, column sep=small, ampersand replacement=\&]
			X \& Z \arrow[rightarrowtail, l] \arrow[d] \\
			\& W \arrow[rightarrowtail, lu]
	\end{tikzcd}}\right\}, \\
	\Obj\left(S_{X/}\right) &:= \left\{\text{morphisms}\;X\rightarrowtail Z\right\}, & \Mor(S_{X/}) &:= \left\{\text{commutative diagrams}\;
	% segment-id: o014.li2.chapter1.gabriel-zisman-localization.g007
	{\begin{tikzcd}[row sep=small, column sep=small, ampersand replacement=\&]
			X \arrow[rightarrowtail, r] \arrow[rightarrowtail, rd] \& Z \arrow[d] \\
			\& W
	\end{tikzcd}}\right\}.
\end{align*}
Note that if $S$ is a left multiplicative system, then
$\left(S^{\opp}\right)_{X/}=(S_{/X})^{\opp}$.

% segment-id: o014.li2.chapter1.gabriel-zisman-localization.le001
\begin{lemma}\label{prop:S-filtered}
	If $S$ is a left (or right) multiplicative system, then
	$S_{/X}^{\opp}$ (or $S_{X/}$) is a filtered category.
\end{lemma}
% segment-id: o014.li2.chapter1.gabriel-zisman-localization.pf004
\begin{proof}
	By the duality observed above, we need only discuss the case of a right
	multiplicative system. For any objects $i:X\to Z$ and $j:X\to Z'$ of
	$S_{X/}$, use the corresponding version of (S3) to construct a commutative
	diagram
	% segment-id: o014.li2.chapter1.gabriel-zisman-localization.g008
	\[\begin{tikzcd}
		Z \arrow[r, "{j'}"] & W \\
		X \arrow[rightarrowtail, u, "i"] \arrow[rightarrowtail, r, "j"'] & Z' \arrow[rightarrowtail, u, "{i'}"']
	\end{tikzcd}\]
	Then $k:=i'j\in S$ gives an object of $S_{X/}$, together with morphisms
	$i\xrightarrow{j'}k\xleftarrow{i'}j$ in $S_{X/}$.

	Let $i,j$ be as above. For any pair of morphisms in $S_{X/}$,
	% segment-id: o014.li2.chapter1.gabriel-zisman-localization.g009
	$\begin{tikzcd}[row sep=small]
		Z \arrow[shift left, rr, "f"] \arrow[shift right, rr, "g"'] & & Z' \\
		& X \arrow[lu, "i"] \arrow[ru, "j"'] &
	\end{tikzcd}$
	(a commutative diagram), the corresponding version of (S4) gives
	$h:Z'\rightarrowtail W$ such that $hf=hg$. Put
	$k:=hj:X\rightarrowtail W$ and regard this as an object of $S_{X/}$. We
	then obtain the commutative diagram in $S_{X/}$
	% segment-id: o014.li2.chapter1.gabriel-zisman-localization.g010
	$\begin{tikzcd}[column sep=small]
		i \arrow[shift left, r, "f"] \arrow[shift right, r, "g"'] & j \arrow[r, "h"] & k
	\end{tikzcd}$.
	These are exactly the defining conditions for a filtered category.
\end{proof}

% segment-id: o014.li2.chapter1.gabriel-zisman-localization.p008
Let $X,Y\in\Obj(\mathcal{C})$. We now define sets $M_{X,Y}^l$ and
$M_{X,Y}^r$ for left and right multiplicative systems, respectively. Note
that they need not be small sets.

% segment-id: o014.li2.chapter1.gabriel-zisman-localization.l003
\begin{itemize}
	% segment-id: o014.li2.chapter1.gabriel-zisman-localization.i007
	\item Let $S$ be a left multiplicative system. Abbreviate a diagram in
	$\mathcal{C}$ of the form
	% segment-id: o014.li2.chapter1.gabriel-zisman-localization.g011
	\[\begin{tikzcd} X & Z \arrow[rightarrowtail, l, "s"'] \arrow[r, "a"] & Y \end{tikzcd}\]
	by $(Z;s,a)$. These data form the set $M_{X,Y}^l$.
	% segment-id: o014.li2.chapter1.gabriel-zisman-localization.i008
	\item Let $S$ be a right multiplicative system. Abbreviate a diagram in
	$\mathcal{C}$ of the form
	% segment-id: o014.li2.chapter1.gabriel-zisman-localization.g012
	\[\begin{tikzcd} X \arrow[r, "a"] & Z & Y \arrow[rightarrowtail, l, "s"'] \end{tikzcd}\]
	by $(Z;a,s)$. These data form the set $M_{X,Y}^r$.
\end{itemize}

% segment-id: o014.li2.chapter1.gabriel-zisman-localization.p009
Define a binary relation $\sim$ as follows:
$(Z;s,a)\sim(Z';s',a')$ (or $(Z;a,s)\sim(Z';a',s')$) if and only if there
is a commutative diagram of the form
% segment-id: o014.li2.chapter1.gabriel-zisman-localization.q004
\begin{equation}\label{eqn:localization-sim}
	\begin{array}{cc}
		% segment-id: o014.li2.chapter1.gabriel-zisman-localization.g013
		\begin{tikzcd}
			& Z \arrow[rightarrowtail, ld, "s"'] \arrow[rd, "a"] & \\
			X & W \arrow[rightarrowtail, l] \arrow[r] \arrow[u] \arrow[d] & Y \\
			& Z' \arrow[rightarrowtail, lu, "{s'}"] \arrow[ru, "{a'}"'] &
		\end{tikzcd} &
		% segment-id: o014.li2.chapter1.gabriel-zisman-localization.g014
		\begin{tikzcd}
			& Z \arrow[d] & \\
			X \arrow[r] \arrow[ru, "a"] \arrow[rd, "{a'}"'] & W & Y \arrow[rightarrowtail, l] \arrow[rightarrowtail, lu, "s"'] \arrow[rightarrowtail, ld, "{s'}"] \\
			& Z' \arrow[u] &
		\end{tikzcd} \\
		\text{(left multiplicative system)} & \text{(right multiplicative system)}
	\end{array}
\end{equation}

% segment-id: o014.li2.chapter1.gabriel-zisman-localization.p010
For every $f:X\to Y$, condition (S1) implies
$(X;\identity_X,f)\in M_{X,Y}^l$ and
$(X;f,\identity_Y)\in M_{X,Y}^r$. The two definitions are plainly dual.

% segment-id: o014.li2.chapter1.gabriel-zisman-localization.le002
\begin{lemma}\label{prop:MXY-equiv}
	For a left (or right) multiplicative system $S$ and arbitrary $X,Y$, the
	relation $\sim$ defined above is an equivalence relation on $M_{X,Y}^l$
	(or $M_{X,Y}^r$). In fact, there are bijections
	% segment-id: o014.li2.chapter1.gabriel-zisman-localization.g015
	\[\begin{tikzcd}[row sep=tiny, column sep=small]
		M_{X, Y}^l /\sim \arrow[r, "1:1"] & \displaystyle\varinjlim_{\substack{[X \leftarrowtail Z] \\ \in \Obj(S_{/X}^{\opp})}} \Hom(Z, Y) , \\
		{[Z; s, a]} \arrow[mapsto, r] & {\left[ Z \xrightarrow{a} Y \right]}
	\end{tikzcd} \quad
	% segment-id: o014.li2.chapter1.gabriel-zisman-localization.g016
	\begin{tikzcd}[row sep=tiny, column sep=small]
		M_{X, Y}^r /\sim \arrow[r, "1:1"] & \displaystyle\varinjlim_{\substack{[Y \rightarrowtail Z] \\ \in \Obj(S_{Y/})}} \Hom(X, Z) \\
		{[Z; a, s]} \arrow[mapsto, r] & {\left[ X \xrightarrow{a} Z \right]}
	\end{tikzcd}\]
	where $[Z;s,a]$ (or $[Z;a,s]$) denotes the equivalence class containing
	$(Z;s,a)$ (or $(Z;a,s)$).
\end{lemma}
% segment-id: o014.li2.chapter1.gabriel-zisman-localization.pf005
\begin{proof}
	Consider the functor $\alpha:S_{/X}^{\opp}\to\cate{Set}$ sending an object
	$X\leftarrowtail Z$ to $\Hom(Z,Y)$. By definition,
	% segment-id: o014.li2.chapter1.gabriel-zisman-localization.q005
	\[ M_{X, Y}^l = \bigsqcup_{X \leftarrowtail Z} \alpha(X \leftarrowtail Z), \]
	and since $S_{/X}^{\opp}$ is filtered (Lemma \ref{prop:S-filtered}), it is
	easy to check that the binary relation $\sim$ on $M_{X,Y}^l$ arising from
	\eqref{eqn:localization-sim} translates into the equivalence relation in
	Proposition \ref{prop:filtered-union}. This also proves the bijection with
	$\varinjlim$. The argument for the $M_{X,Y}^r$ version is identical.

	When discussing filtered colimits of sets, \S\ref{sec:filtered-indlim}
	assumes that the indexing category is small, so that the resulting colimit
	still lies in $\cate{Set}$. Here $S_{/X}$ and $S_{Y/}$ need not be small
	categories. The assertion itself, however, does not depend on set size, and
	the argument can be formulated so as not to depend on a choice of
	Grothendieck universe.
\end{proof}

% segment-id: o014.li2.chapter1.gabriel-zisman-localization.p011
The equivalence class $[Z;s,a]$ (or $[Z;a,s]$) should be understood as the
morphism $as^{-1}$ (or $s^{-1}a$) in the category $\mathcal{C}[S^{-1}]$ to be
constructed.

% segment-id: o014.li2.chapter1.gabriel-zisman-localization.dp001
\begin{definition-proposition}\label{prop:localized-composition}
	Let $S$ be a left multiplicative system, let $X,Y$ be arbitrary, and take
	$(U;s,a)\in M_{X,Y}^l$ and $(V;t,b)\in M_{Y,Z}^l$. Use (S3) to extend them
	to the following commutative diagram:
	% segment-id: o014.li2.chapter1.gabriel-zisman-localization.g017
	\[\begin{tikzcd}[row sep=small]
		& & W \arrow[rightarrowtail, ld, "r"'] \arrow[rd, "c"] & & \\
		& U \arrow[rightarrowtail, ld, "s"'] \arrow[rd, "a"'] & & V \arrow[rightarrowtail, ld, "t"] \arrow[rd, "b"] & \\
		X & & Y & & Z
	\end{tikzcd}\]
	Then
	$[V;t,b]\circ[U;s,a]:=[W;sr,bc]\in M_{X,Z}^l/\sim$ depends only on
	$[U;s,a]$ and $[V;t,b]$. The dual statement holds for right multiplicative
	systems.
\end{definition-proposition}
% segment-id: o014.li2.chapter1.gabriel-zisman-localization.pf006
\begin{proof}
	First fix $(V;t,b)$ and vary $(U;s,a)$ within its equivalence class. In
	view of \eqref{eqn:localization-sim}, it suffices, for a commutative diagram
	% segment-id: o014.li2.chapter1.gabriel-zisman-localization.g018
	\[\begin{tikzcd}
		& U \arrow[rightarrowtail, ld, "s"'] \arrow[rd, "a"] & W \arrow[rightarrowtail, l, "r"'] \arrow[rd, "c"] & & \\
		X & & Y & V \arrow[r, "b"] \arrow[rightarrowtail, l, "t"'] & Z . \\
		& U' \arrow[rightarrowtail, lu, "{s'}"] \arrow[ru, "{a'}"] \arrow[uu, "x"] & W' \arrow[rightarrowtail, l, "{r'}"] \arrow[ru, "{c'}"'] & &
	\end{tikzcd}\]
	to prove $(W;sr,bc)\sim(W';s'r',bc')$. Apply (S3) to obtain an object
	$Q$ and morphisms $q,k$ such that
	% segment-id: o014.li2.chapter1.gabriel-zisman-localization.g019
	\[\begin{tikzcd}
		Q \arrow[rightarrowtail, d, "k"'] \arrow[r, "q"] & W \arrow[rightarrowtail, d, "r"] \\
		W' \arrow[r, "{xr'}"'] & U
	\end{tikzcd} \quad \text{commutes}. \]

	Combining this with the preceding diagram gives
	$tc'k=a'r'k=axr'k=arq=tcq$. By (S4), there is a
	$w:P\rightarrowtail Q$ such that $c'kw=cqw$. It is then easy to verify that
	% segment-id: o014.li2.chapter1.gabriel-zisman-localization.g020
	\[\begin{tikzcd}[column sep=large, row sep=large]
		& W \arrow[rightarrowtail, ld, "sr"'] \arrow[rd, "bc"] & \\
		X & P \arrow[r, "{bc'kw}" description] \arrow[rightarrowtail, l, "{s'r'kw}" description] \arrow[u, "qw" description] \arrow[d, "{kw}" description] & Z \\
		& W' \arrow[rightarrowtail, lu, "{s'r'}"] \arrow[ru, "{bc'}"'] &
	\end{tikzcd} \quad \text{commutes, that is,}\; (W;sr,bc)\sim(W';s'r',bc'). \]

	The case in which $(U;s,a)$ is fixed and $(V;t,b)$ varies within its
	equivalence class is similar.
\end{proof}

% segment-id: o014.li2.chapter1.gabriel-zisman-localization.dt001
\begin{definition-theorem}\label{def:fraction-localization}
	Let $S$ be a left multiplicative system in a category $\mathcal{C}$.
	% segment-id: o014.li2.chapter1.gabriel-zisman-localization.l004
	\begin{enumerate}[(i)]
		% segment-id: o014.li2.chapter1.gabriel-zisman-localization.i009
		\item There is a category $\mathcal{C}[S^{-1}]^l$ (possibly a large
		category) whose object set is $\Obj(\mathcal{C})$ and whose $\Hom$ set
		between two objects $X,Y$ is $M_{X,Y}^l/\sim$. The identity morphism of
		$X$ is $[X;\identity_X,\identity_X]$, and composition is given by the
		binary operation in Definition--Proposition
		\ref{prop:localized-composition}.
		% segment-id: o014.li2.chapter1.gabriel-zisman-localization.i010
		\item There is a functor
		$Q^l:\mathcal{C}\to\mathcal{C}[S^{-1}]^l$ that is the identity on
		objects and sends a morphism $f:X\to Y$ to
		$[X;\identity_X,f]$.
	\end{enumerate}
	The dual statement holds for a right multiplicative system $S$; the
	corresponding data are denoted
	$Q^r:\mathcal{C}\to\mathcal{C}[S^{-1}]^r$.
\end{definition-theorem}
% segment-id: o014.li2.chapter1.gabriel-zisman-localization.pf007
\begin{proof}
	By duality, we need only consider a left multiplicative system.

	For (i), it suffices to verify the properties of composition in
	$\mathcal{C}[S^{-1}]^l$. The operation in Definition--Proposition
	\ref{prop:localized-composition} plainly makes
	$[X;\identity_X,\identity_X]$ an identity morphism, so only associativity
	remains. Take $(A;s,a)\in M_{X,Y}^l$, $(B;t,b)\in M_{Y,Z}^l$, and
	$(C;u,c)\in M_{Z,W}^l$. Applying (S3) three times gives the commutative
	diagram
	% segment-id: o014.li2.chapter1.gabriel-zisman-localization.g021
	\[\begin{tikzcd}[row sep=small]
		& & & \bullet \arrow[rightarrowtail, ld] \arrow[rd] & & & \\
		& & \bullet \arrow[rightarrowtail, ld] \arrow[rd] & & \bullet \arrow[rightarrowtail, ld] \arrow[rd] & & \\
		& A \arrow[rightarrowtail, ld, "s"'] \arrow[rd, "a"] & & B \arrow[rightarrowtail, ld, "t"'] \arrow[rd, "b"] & & C \arrow[rightarrowtail, ld, "u"'] \arrow[rd, "c"] & \\
		X & & Y & & Z & & W
	\end{tikzcd}\]
	The whole diagram gives an element of $M_{X,W}^l$; its lower-left part
	gives a representative of $[B;t,b]\circ[A;s,a]$, and its lower-right part
	gives a representative of $[C;u,c]\circ[B;t,b]$. This proves
	associativity:
	% segment-id: o014.li2.chapter1.gabriel-zisman-localization.q006
	\[ [C;u,c]\circ([B;t,b]\circ[A;s,a])=([C;u,c]\circ[B;t,b])\circ[A;s,a]. \]

	For (ii), it remains to check that $Q^l$ respects identities and composition.
	It preserves identities because $[X;\identity_X,\identity_X]$ has the
	required identity property. Now consider morphisms
	$X\xrightarrow{f}Y\xrightarrow{g}Z$ in
	$\mathcal{C}$. The diagram
	% segment-id: o014.li2.chapter1.gabriel-zisman-localization.g022
	\[\begin{tikzcd}[row sep=small]
		& & X \arrow[equal, ld] \arrow[rd, "f"] & & \\
		& X \arrow[equal, ld] \arrow[rd, "f"] & & Y \arrow[equal, ld] \arrow[rd, "g"] & \\
		X & & Y & & Z
	\end{tikzcd}\]
	shows that $Q^l(gf)=Q^l(g)Q^l(f)$.
\end{proof}

% segment-id: o014.li2.chapter1.gabriel-zisman-localization.th001
\begin{theorem}[P.\ Gabriel, M.\ Zisman]\label{prop:Gabriel-Zisman}\index{localization}
	Let $S$ be a left (or right) multiplicative system in $\mathcal{C}$. Then
	$\left(\mathcal{C}[S^{-1}]^l,Q^l\right)$ (or
	$\left(\mathcal{C}[S^{-1}]^r,Q^r\right)$) is a localization of
	$\mathcal{C}$ at $S$.
\end{theorem}
% segment-id: o014.li2.chapter1.gabriel-zisman-localization.pf008
\begin{proof}
	By duality, we only verify the conditions of Definition
	\ref{def:cat-localization} for a left multiplicative system. Write
	$Q=Q^l$. If $f:X\to Y$ belongs to $S$, the diagram
	% segment-id: o014.li2.chapter1.gabriel-zisman-localization.g023
	\[ \begin{tikzcd}
		& X \arrow[rightarrowtail, ld, "f"'] \arrow[rd, "f"] \arrow[d, "f" description] & \\
		Y & Y \arrow[equal, l] \arrow[equal, r] & Y
	\end{tikzcd}\]
	shows that $[X;f,f]=[Y;\identity_Y,\identity_Y]$. Hence the following
	diagrams show that $[X;f,\identity_X]=[X;\identity_X,f]^{-1}$:
	% segment-id: o014.li2.chapter1.gabriel-zisman-localization.g024
	\[\begin{tikzcd}[column sep=small, row sep=small]
		& & X \arrow[equal, ld] \arrow[equal, rd] & & \\
		& X \arrow[equal, ld] \arrow[rd, "f"] & & X \arrow[rightarrowtail, ld, "f"'] \arrow[equal, rd] & \\
		X & & Y & & X
	\end{tikzcd} \quad
	% segment-id: o014.li2.chapter1.gabriel-zisman-localization.g025
	\begin{tikzcd}[column sep=small, row sep=small]
		& & X \arrow[equal, ld] \arrow[equal, rd] & & \\
		& X \arrow[rightarrowtail, ld, "f"'] \arrow[equal, rd] & & X \arrow[equal, ld] \arrow[rightarrowtail, rd, "f"] & \\
		Y & & X & & Y
	\end{tikzcd}\]
	Thus $Q$ does send the elements of $S$ to isomorphisms.

	Now let $F:\mathcal{C}\to\mathcal{D}$ be a functor sending $S$ to
	isomorphisms. Define
	$F[S^{-1}]:\mathcal{C}[S^{-1}]^l\to\mathcal{D}$ as follows: on objects it
	sends $X$ to $FX$, and on morphisms it sends
	$[U;s,a]\in\Hom_{\mathcal{C}[S^{-1}]^l}(X,Y)$ to
	$(Fa)(Fs)^{-1}\in\Hom_{\mathcal{D}}(FX,FY)$\footnote{Translator's note:
	the source has $\Hom_{\mathcal{D}}(X,Y)$. Since
	$(Fa)(Fs)^{-1}:FX\to FY$, the source and target objects have been corrected
	to $FX$ and $FY$ here (O014-C018).}. Clearly, $F[S^{-1}]$ sends identity
	morphisms in $\mathcal{C}[S^{-1}]^l$ to identity morphisms in $\mathcal{D}$.
	Apply $F$ to the diagram in Definition--Proposition
	\ref{prop:localized-composition} and invert the arrows in $F(S)$; it follows
	at once that $F[S^{-1}]$ preserves composition. The equality
	$F[S^{-1}]Q=F$ is also clear. Since every morphism in
	$\mathcal{C}[S^{-1}]^l$ has the form $(Qa)(Qs)^{-1}$, no other definition
	of $F[S^{-1}]$ is possible.
\end{proof}

% segment-id: o014.li2.chapter1.gabriel-zisman-localization.c002
\begin{convention}
	If $S$ is a multiplicative system, Proposition
	\ref{prop:cat-localization-unique} shows that the two localizations can be
	identified; we write the resulting functor as
	$Q:\mathcal{C}\to\mathcal{C}[S^{-1}]$.
\end{convention}

% segment-id: o014.li2.chapter1.gabriel-zisman-localization.co001
\begin{corollary}\label{prop:fraction-zero}
	Let $S$ be a left (or right) multiplicative system in $\mathcal{C}$.
	Morphisms $f,g:X\to Y$ in $\mathcal{C}$ satisfy $Q^lf=Q^lg$ (or
	$Q^rf=Q^rg$) if and only if there exists $s\in S$ such that $fs=gs$ (or
	$sf=sg$).
\end{corollary}
% segment-id: o014.li2.chapter1.gabriel-zisman-localization.pf009
\begin{proof}
	It suffices to check the ``only if'' direction for a left multiplicative
	system $S$. If $Q^lf=Q^lg$, the construction of
	$\mathcal{C}[S^{-1}]^l$ gives a commutative diagram
	% segment-id: o014.li2.chapter1.gabriel-zisman-localization.g026
	\[\begin{tikzcd}
		& X \arrow[equal, ld] \arrow[rd, "f"] & \\
		X & U \arrow[rightarrowtail, l, "s"'] \arrow[r, "a"] \arrow[d] \arrow[u] & Y \\
		& X \arrow[equal, lu] \arrow[ru, "g"'] &
	\end{tikzcd}\]
	from which $fs=a=gs$.
\end{proof}

% segment-id: o014.li2.chapter1.gabriel-zisman-localization.co002
\begin{corollary}\label{prop:fraction-mono-epi}
	Let $S$ be a left (or right) multiplicative system in $\mathcal{C}$. Then
	$Q^l$ (or $Q^r$) sends monomorphisms to monomorphisms (or epimorphisms to
	epimorphisms).
\end{corollary}
% segment-id: o014.li2.chapter1.gabriel-zisman-localization.pf010
\begin{proof}
	We prove only the case of a left multiplicative system. Let $f:X\to Y$ be
	a monomorphism in $\mathcal{C}$. Take
	$\alpha,\beta:Q^lW\rightrightarrows Q^lX$ such that
	$(Q^lf)\alpha=(Q^lf)\beta$. Conditions (S2) and (S3) allow us to represent
	$\alpha$ and $\beta$ by elements $(U;s,a)$ and $(U;s,b)$, respectively,
	of $M^l_{W,X}$; the details are left to the reader as an exercise. The
	given equality then implies $Q^l(fa)=Q^l(fb)$. By Corollary
	\ref{prop:fraction-zero}, there is a $t\in S$ such that $fat=fbt$. Hence
	$at=bt$, and therefore $\alpha=\beta$.
\end{proof}

% segment-id: o014.li2.chapter1.gabriel-zisman-localization.r001
\begin{remark}\label{rem:localization-size}
	\index{size issues}
	In the basic operations of category theory, we want to avoid large
	categories whenever possible. Definition--Theorem
	\ref{def:fraction-localization} says that $\mathcal{C}[S^{-1}]^l$ (or
	$\mathcal{C}[S^{-1}]^r$) may be a large category. The reason is that the
	$\varinjlim$ in Lemma \ref{prop:MXY-equiv} need not be small, so its value
	need not lie in $\cate{Set}$. This is a somewhat delicate issue.

	If $S$ satisfies the following condition, however, then
	$\mathcal{C}[S^{-1}]^l$ (or $\mathcal{C}[S^{-1}]^r$) is not large: assume
	that $S_{/X}^{\opp}$ (or $S_{X/}$) has a small cofinal subcategory
	(Definition \ref{def:cofinal}) for every $X\in\Obj(\mathcal{C})$. By
	Proposition \ref{prop:cofinal-lim}, the colimit in question may be
	restricted to that subcategory, yielding an object of $\cate{Set}$. If
	$\mathcal{C}$ itself is small, this condition holds automatically.
\end{remark}

% segment-id: o014.li2.chapter1.gabriel-zisman-localization.le003
\begin{lemma}\label{prop:localization-prod}
	Let $S$ be a left (or right) multiplicative system in $\mathcal{C}$. Then
	the localization functor $Q^l$ (or $Q^r$) preserves finite $\varprojlim$
	(or finite $\varinjlim$).
\end{lemma}
% segment-id: o014.li2.chapter1.gabriel-zisman-localization.pf011
\begin{proof}
	We discuss only the case of a left multiplicative system. Let $J$ be a
	finite category, and suppose that the limit associated with a functor
	$\beta:J^{\opp}\to\mathcal{C}$ exists. By Lemma \ref{prop:MXY-equiv}, for
	every $X\in\Obj(\mathcal{C})$ there are canonical bijections
	% segment-id: o014.li2.chapter1.gabriel-zisman-localization.q007
	\begin{multline*}
		\Hom_{\mathcal{C}[S^{-1}]^l}\left(Q^lX,Q^l\varprojlim\beta\right)
		\simeq \varinjlim_{X\leftarrowtail Z}\Hom_{\mathcal{C}}\left(Z,\varprojlim\beta\right) \\
		\simeq \varinjlim_{X\leftarrowtail Z}\varprojlim_{j\in\Obj(J)}
		\Hom_{\mathcal{C}}\left(Z,\beta(j)\right).
	\end{multline*}

	Applying Lemma \ref{prop:S-filtered} and Proposition
	\ref{prop:filtered-finite-exchange} to the last term transforms it further
	into
	% segment-id: o014.li2.chapter1.gabriel-zisman-localization.q008
	\[ \varprojlim_{j \in \Obj(J)} \varinjlim_{X \leftarrowtail Z} \Hom_{\mathcal{C}}\left( Z, \beta(j) \right) \simeq \varprojlim_{j \in \Obj(J)} \Hom_{\mathcal{C}[S^{-1}]^l} \left(Q^l X, Q^l \beta(j) \right). \]
	Strictly speaking, \S\ref{sec:filtered-indlim} works in $\cate{Set}$,
	whereas the $\Hom$ sets here need not lie in $\cate{Set}$. This is only a
	minor technical issue; see the related discussion in the proof of Lemma
	\ref{prop:MXY-equiv}.
\end{proof}

% segment-id: o014.li2.chapter1.gabriel-zisman-localization.th002
\begin{theorem}\label{prop:localization-additivity}
	Let $\mathcal{C}$ be an $\cate{Ab}$-category and let $S$ be a left (or
	right) multiplicative system in it.
	% segment-id: o014.li2.chapter1.gabriel-zisman-localization.l005
	\begin{enumerate}[(i)]
		% segment-id: o014.li2.chapter1.gabriel-zisman-localization.i011
		\item The category $\mathcal{C}[S^{-1}]^l$ (or
		$\mathcal{C}[S^{-1}]^r$) has a canonical $\cate{Ab}$-category structure
		making $Q^l$ (or $Q^r$) an additive, essentially surjective functor.
		% segment-id: o014.li2.chapter1.gabriel-zisman-localization.i012
		\item If $\mathcal{C}$ is an additive category and $S$ is a
		multiplicative system, then $\mathcal{C}[S^{-1}]$ is also an additive
		category.
		% segment-id: o014.li2.chapter1.gabriel-zisman-localization.i013
		\item The preceding statements remain true if $\cate{Ab}$-categories are
		replaced by $\Bbbk$-linear categories, where $\Bbbk$ is a commutative
		ring. If $F$ in the universal property of Definition
		\ref{def:cat-localization} is $\Bbbk$-linear, then so is the associated
		functor $F[S^{-1}]$.
	\end{enumerate}
\end{theorem}
% segment-id: o014.li2.chapter1.gabriel-zisman-localization.pf012
\begin{proof}
	For (i), it suffices to treat a left multiplicative system. We must define
	addition on the morphism sets of $\mathcal{C}[S^{-1}]$. Let
	\[
		f,g\in\Hom_{\mathcal{C}[S^{-1}]^l}(X,Y).
	\]
	Since $S_{/X}^{\opp}$ is filtered (Lemma \ref{prop:S-filtered}), there are
	a morphism $s:U\rightarrowtail X$ and morphisms $a_1,a_2:U\to Y$ in
	$\mathcal{C}$ such that $f=[U;s,a_1]$ and $g=[U;s,a_2]$. Define
	% segment-id: o014.li2.chapter1.gabriel-zisman-localization.q009
	\[ f+g:=[U;s,a_1+a_2]\in\Hom_{\mathcal{C}[S^{-1}]}(X,Y). \]
	In particular, a zero morphism can be written $[U;s,0]$, and
	$-[U;s,a]=[U;s,-a]$. One must still prove that $f+g$ is independent of the
	choice of representatives; this too follows from the filteredness of
	$S_{/X}^{\opp}$, and the routine details are omitted. Once this operation
	is known to be well defined, the axioms for an $\cate{Ab}$-category and the
	additivity of $Q$ are easy to verify.

	Statement (ii) follows from (i), the definition of an additive category
	(Definition \ref{def:additive-category}), and Lemma
	\ref{prop:localization-prod}.

	For (iii), if $\mathcal{C}$ is $\Bbbk$-linear, define scalar multiplication
	on $\Hom_{\mathcal{C}[S^{-1}]}$ by
	$t\cdot[U;s,a]=[U;s,ta]$ for $t\in\Bbbk$. The required properties are
	immediate.
\end{proof}

% segment-id: o014.li2.chapter1.gabriel-zisman-localization.e001
\begin{example}[Localization of rings]
	Giving a ring $R$ is equivalent to giving an $\cate{Ab}$-category
	$\mathcal{C}$ with a single object $\star$ such that
	$R=\End_{\mathcal{C}}(\star)$. If $R$ is commutative, the conditions that
	$S\subset R=\Mor(\mathcal{C})$ be a left or a right multiplicative system
	coincide. In this case, the ring corresponding to the $\cate{Ab}$-category
	$\mathcal{C}[S^{-1}]$ is precisely the localization $R[S^{-1}]$ of $R$;
	see \cite[\S 5.3]{Li1}. For a general ring $R$, the theory in this section
	provides a construction of noncommutative localizations. The related theory
	is also a branch of noncommutative ring theory; see \cite[\S 10A]{Lam99}
	for a detailed account.
\end{example}

% segment-id: o014.li2.chapter1.gabriel-zisman-localization.e002
\begin{example}[Central localization]\label{eg:central-localization}
	Write the center of an additive category $\mathcal{C}$ as
	$Z(\mathcal{C})$; this is a commutative ring. See the discussion preceding
	Proposition \ref{prop:k-linear-center}. For any multiplicative subset $S$
	of $Z(\mathcal{C})$, it is easy to check that the set of morphisms
	% segment-id: o014.li2.chapter1.gabriel-zisman-localization.q010
	\[ \{s_X \in \End_{\mathcal{C}}(X): s \in S, \; X \in \Obj(\mathcal{C}) \} \]
	is a multiplicative system. The corresponding category
	$\mathcal{C}[S^{-1}]$ can of course be constructed using the calculus of
	fractions above, but a shorter construction is
	% segment-id: o014.li2.chapter1.gabriel-zisman-localization.q011
	\[ \Obj(\mathcal{C}[S^{-1}]) := \Obj(\mathcal{C}), \quad \Hom_{\mathcal{C}[S^{-1}]}(X, Y) := \Hom_{\mathcal{C}}(X, Y) \dotimes{Z(\mathcal{C})} Z(\mathcal{C})[S^{-1}], \]
	with composition defined in the usual way. The relevant checks are left as
	an exercise for this chapter.
\end{example}

% segment-id: o014.li2.chapter1.gabriel-zisman-localization.r002
\begin{remark}\label{rem:reflexive-localization}
	\index{localization!reflective}
	For general $\mathcal{C}$ and $S$, the case in which the localization
	functor $Q:\mathcal{C}\to\mathcal{C}[S^{-1}]$ has a fully faithful right
	adjoint is especially interesting and important. A localization with this
	property is called a \emph{reflective localization}. The
	Gabriel--Popescu theorem to be introduced in Appendix \S~A.3 is a
	typical example.
\end{remark}

% segment-id: o014.li2.chapter1.gabriel-zisman-localization.r003
\begin{remark}[Localization of product categories]\label{rem:localization-product-cat}
	Finally, let $\mathcal{C}_i$ be categories equipped with left (or right)
	multiplicative systems $S_i$, for $i=1,2$, and write their localization
	functors as $Q_i:\mathcal{C}_i\to\mathcal{C}_i[S_i^{-1}]$. Clearly,
	$S_1\times S_2$ is also a left (or right) multiplicative system in
	$\mathcal{C}_1\times\mathcal{C}_2$. From the construction in
	Definition--Theorem \ref{def:fraction-localization}, one readily checks that
	% segment-id: o014.li2.chapter1.gabriel-zisman-localization.q012
	\[ (Q_1, Q_2): \mathcal{C}_1 \times \mathcal{C}_2 \to \mathcal{C}_1[S_1^{-1}] \times \mathcal{C}_2[S_2^{-1}] \]
	is the corresponding localization. This statement generalizes to products
	of arbitrary families of categories.
\end{remark}
